\documentclass[11pt]{article}
\usepackage[margin=1in]{geometry}
\usepackage[T1]{fontenc}
\usepackage[utf8]{inputenc}
\usepackage{amsmath,amssymb,amsthm}
\usepackage{enumitem}

\title{ICPC World Finals 2022\\Q. Doing the Container Shuffle}
\author{}
\date{}

\begin{document}
\maketitle

\section*{Problem Summary}

Containers $1,2,\dots,n$ are placed independently onto one of two stacks. They are later requested in a
fixed order $a_1,\dots,a_n$. If the next requested container is not on top of its stack, we repeatedly move
containers from that stack to the other stack until it is exposed. We must compute the expected total
number of moves, including the final move from a stack to a truck.

\section*{Key Observations}

\begin{itemize}[leftmargin=*]
    \item At any moment, the two stacks can be viewed as one \emph{joint order}:
    first stack from bottom to top, then second stack from top to bottom.
    Moving the top container from one stack to the other preserves this joint order; it only changes where
    the ``cut'' between the two stacks lies.
    \item Before unloading $a_{i+1}$, the containers that must be moved are exactly the still-present
    containers lying in the joint-order interval between $a_i$ and $a_{i+1}$. For $a_1$, the blocking
    containers are exactly those above $a_1$ on its stack.
    \item A container $v$ can lie in that interval only if it was unloaded from the ship after
    $\min(a_i,a_{i+1})$. For such a container, the random choice of stack puts it in the interval with
    probability $1/2$.
    \item Therefore
    \[
    \mathbb{E}[\text{moves before unloading } a_{i+1}]
    = \frac{1}{2}\left|\{v : v \text{ is requested after } a_{i+1},\ v > \min(a_i,a_{i+1})\}\right|.
    \]
    For the first request $a_1$, the same formula holds with just $v > a_1$.
\end{itemize}

\section*{Algorithm}

\begin{enumerate}[leftmargin=*]
    \item Start the answer at $n$, since every requested container must eventually be moved from a stack
    to a truck exactly once.
    \item Process the truck order from right to left, maintaining in a Fenwick tree the set of container
    labels that appear later in the truck order.
    \item For position $i$:
    \begin{itemize}[leftmargin=*]
        \item if $i=1$, use threshold $a_1$;
        \item otherwise use threshold $\min(a_{i-1},a_i)$.
    \end{itemize}
    \item Query how many later-requested containers have label larger than that threshold, divide by $2$,
    and add it to the answer.
    \item Insert $a_i$ into the Fenwick tree and continue.
\end{enumerate}

\section*{Correctness Proof}

We prove that the algorithm returns the correct answer.

\paragraph{Lemma 1.}
After any sequence of stack-to-stack moves, the relative order of the containers in the joint order is the
same as initially, except that some containers may already have been removed to trucks.

\paragraph{Proof.}
Moving the top container of one stack to the top of the other stack removes that container from one end of
the joint order and inserts it at the other end. The sequence of the remaining containers is unchanged.
Removing a requested container simply deletes it from the joint order. Repeating these operations proves
the claim. \qed

\paragraph{Lemma 2.}
Before unloading $a_{i+1}$, the containers that must be moved are exactly the still-present containers in
the initial joint-order interval between $a_i$ and $a_{i+1}$.

\paragraph{Proof.}
By Lemma 1, after $a_i$ has been removed, the still-present containers keep the same relative order as in
the initial joint order with already removed containers deleted. To expose $a_{i+1}$, we must move exactly
the containers lying between the current cut and $a_{i+1}$, which is the same as the initial interval
between $a_i$ and $a_{i+1}$ after deleting already removed containers. \qed

\paragraph{Lemma 3.}
Fix a later-requested container $v$. It belongs to the interval of Lemma 2 with probability $1/2$ if
$v > \min(a_i,a_{i+1})$, and with probability $0$ otherwise.

\paragraph{Proof.}
If $v < \min(a_i,a_{i+1})$, then $v$ was loaded onto a stack before both endpoints, so it lies below both
of them on whichever stack contains it and cannot be between them in the joint order.

If $v > \min(a_i,a_{i+1})$, then $v$ was loaded later than at least one endpoint. Its random stack choice is
independent of all others, and exactly one of the two stack choices places it in the interval between the
two endpoints in the initial joint order. Hence the probability is $1/2$. \qed

\paragraph{Theorem.}
The algorithm outputs the correct expected total number of moves.

\paragraph{Proof.}
By Lemma 2, the number of stack-to-stack moves before unloading $a_{i+1}$ equals the number of
still-present containers in the corresponding interval. By Lemma 3, the expectation of that number is
exactly one half of the number of later-requested containers whose labels exceed the threshold
$\min(a_i,a_{i+1})$ (or just $a_1$ for the first request). Summing these expectations over all requests and
adding the $n$ inevitable truck moves yields exactly the formula computed by the algorithm. \qed

\section*{Complexity Analysis}

Each of the $n$ positions performs one Fenwick query and one update. The running time is
$O(n \log n)$ and the memory usage is $O(n)$.

\section*{Implementation Notes}

\begin{itemize}[leftmargin=*]
    \item The answer is always a multiple of $0.5$, but the implementation prints it as a floating-point
    number with three digits after the decimal point.
\end{itemize}

\end{document}
